\subsection{Scelta degli sperimentatori}
Per lo svolgimento degli esperimenti, si è deciso di coinvolgere attivamente tutti e tre i componenti del gruppo, suddividendosi il campione di utenti. Al fine di non intimidire il partecipante e ridurre al minimo l'ansia da prestazione, ogni singola sessione di test è stata condotta da un solo sperimentatore, il quale ha ricoperto contemporaneamente il ruolo di osservatore e di facilitatore. 

L'obiettivo principale era mettere a proprio agio l'utente tester, lasciandogli completa libertà di esplorare il sistema e di muoversi in autonomia. Durante le sessioni, gli interventi come facilitatori sono stati ridotti al minimo: ci si è limitati a fornire indicazioni puramente generali solo qualora fosse strettamente necessario o esplicitamente richiesto dall'utente. In questi rari casi, è stata prestata la massima attenzione a non influenzare in alcun modo le sue scelte. In particolare, quando il ruolo di sperimentatore era ricoperto dal progettista originale del sistema, è stata accuratamente evitata qualsiasi spiegazione di natura tecnica che potesse alterare la spontaneità dell'interazione o mascherare un potenziale problema di usabilità. 

Infine, per garantire l'omogeneità e l'affidabilità dei dati raccolti in sessioni parallele, non si sono verificate disparità all'interno del team. Sebbene due dei membri del gruppo non conoscessero l'applicativo in partenza, le fasi precedenti del progetto hanno garantito un'adeguata preparazione: l'esecuzione della valutazione euristica e la successiva partecipazione attiva al processo di riprogettazione hanno permesso di studiare a fondo l'interfaccia. In questo modo, tutti gli sperimentatori hanno acquisito un'esperienza solida e completa sull'applicazione, risultando ugualmente qualificati per guidare e osservare gli utenti durante le prove.